% !TEX root = ../DoAn.tex
\documentclass[../DoAn.tex]{subfiles}
\begin{document}

Nối tiếp bối cảnh thực tiễn ở Chương 1, Chương 2 tập trung khảo sát hiện trạng các công cụ phân tích lưu lượng mạng hiện có để chỉ ra hạn chế và xác định các yêu cầu chức năng cần thiết. Thông qua việc phân tích ca sử dụng và đặc tả các yêu cầu chức năng cũng như phi chức năng, chương này sẽ định hình khung nghiệp vụ cốt lõi, làm cơ sở để lựa chọn công nghệ ở Chương 3 và triển khai thiết kế tại Chương 4.

\section{Khảo sát hiện trạng}
Để xây dựng một giải pháp thiết thực và phù hợp, quá trình khảo sát hiện trạng và xác định yêu cầu hệ thống được thực hiện tập trung vào việc nghiên cứu, phân tích và đánh giá các giải pháp công nghệ cũng như phần mềm phân tích lưu lượng mạng hiện có trên thị trường.

\subsection{Phân tích đặc điểm các giải pháp lưu lượng mạng hiện có}
Hiện nay, các công cụ phục vụ giám sát và phân tích mạng khá đa dạng, được thiết kế tối ưu cho từng mục đích sử dụng cụ thể. 

Đối với các phần mềm bắt và giải mã gói tin trực quan, ưu điểm vượt trội là khả năng bóc tách sâu từng byte dữ liệu thô của rất nhiều giao thức mạng. Tuy nhiên, các công cụ này thường chỉ tập trung hiển thị chi tiết gói tin mà thiếu khả năng tổng hợp trực quan hóa luồng giao tiếp hoặc biểu diễn mối quan hệ giữa các thiết bị dưới dạng topo kết nối mạng. Điều này khiến người học gặp nhiều khó khăn trong việc xâu chuỗi thông tin để có một góc nhìn toàn cảnh.

Đối với các hệ thống giám sát an ninh mạng chuyên dụng, thế mạnh nằm ở việc gom cụm luồng lưu lượng và tự động phân tích hành vi bất thường quy mô lớn. Tuy nhiên, các hệ thống này thường rất nặng nề, yêu cầu kỹ năng cấu hình phức tạp và đòi hỏi tài nguyên phần cứng lớn, không phù hợp để triển khai linh hoạt trên các máy tính cá nhân phục vụ thực hành của học viên.

\subsection{Đánh giá và so sánh các giải pháp hiện tại}
Dựa trên kết quả phân tích các giải pháp trên, các công cụ phân tích lưu lượng mạng phổ biến hiện có được đặt lên bàn cân so sánh với giải pháp đề xuất của đồ án theo các tiêu chí quan trọng tại Bảng \ref{tab:so_sanh_cong_cu}.

\begin{table}[H]
\centering
\caption{Bảng so sánh tính năng giữa các giải pháp hiện có và đồ án đề xuất}
\label{tab:so_sanh_cong_cu}
\resizebox{\textwidth}{!}{%
\begin{tabular}{|p{4.5cm}|c|c|c|}
\hline
\textbf{Tiêu chí so sánh} & \textbf{Công cụ đồ họa gói tin} & \textbf{Hệ thống giám sát luồng} & \textbf{Giải pháp đề xuất} \\ \hline
Bắt và giải mã gói tin & Rất tốt & Hạn chế & Tốt \\ \hline
Thống kê luồng lưu lượng & Khá & Tốt & Tốt \\ \hline
Trực quan hóa topo và dòng dữ liệu & Hạn chế & Khá & Rất tốt \\ \hline
Trích xuất luật tường lửa ACL & Không hỗ trợ & Không hỗ trợ & Tốt \\ \hline
Tích hợp học sâu hỗ trợ phân tích & Không hỗ trợ & Khá & Tốt \\ \hline
Mức độ dễ tiếp cận cho người học & Trung bình & Thấp & Cao \\ \hline
\end{tabular}%
}
\end{table}

Qua bảng so sánh, các giải pháp hiện tại tuy rất mạnh ở các mảng chuyên biệt nhưng lại thiếu tính tích hợp phục vụ nhu cầu đào tạo. Các công cụ đồ họa gói tin chỉ tập trung phân tích byte thô mà thiếu tính năng thống kê luồng hay trực quan hóa topo; trong khi đó, các hệ thống giám sát luồng lại quá nặng nề, phức tạp và khó triển khai cho mục đích học tập cá nhân.

\subsection{Xác định các tính năng phần mềm cần phát triển}
Từ kết quả phân tích các giải pháp hiện tại, có thể thấy hạn chế lớn nhất nằm ở sự rời rạc trong quy trình giám sát mạng, khoảng cách lớn giữa dữ liệu gói tin thô tầng thấp và góc nhìn tổng quan mức luồng, cùng với việc thiếu hụt các công cụ trực quan hóa kết nối và hỗ trợ phân tích thông minh dễ tiếp cận cho học viên.

Do đó, để giải quyết triệt để các hạn chế đã nêu, hệ thống cần tập trung phát triển các khối chức năng thu thập và phân tích lưu lượng. Trước hết, phần mềm phải xây dựng cơ chế thu thập gói tin thời gian thực trực tiếp từ các card mạng vật lý kết hợp khả năng đọc hoặc ghi tệp dữ liệu chuẩn. Bên cạnh đó, quy trình giải mã và tìm kiếm cần được thiết kế linh hoạt, hỗ trợ bóc tách sâu cấu trúc các trường thông tin của nhiều giao thức thông dụng, hiển thị mã hex trực quan và cung cấp bộ lọc cú pháp thông minh phục vụ tra cứu dữ liệu.

Một yêu cầu quan trọng khác là việc xây dựng lớp trực quan hóa trực diện và quản lý luồng lưu lượng, trong đó các gói tin đơn lẻ được tự động nhóm thành các luồng hai chiều để tính toán các tham số thống kê về thời gian, dung lượng và tốc độ. Dữ liệu này sau đó được biểu diễn sinh động qua các sơ đồ tiến trình dòng chảy dữ liệu, sơ đồ mạng liên kết các thiết bị mạng cùng hệ thống biểu đồ thống kê dạng đường, cột, tròn trên giao diện hiển thị, qua đó giúp người học có được góc nhìn toàn cảnh về hoạt động của mạng.

Cuối cùng, phần mềm tích hợp các tính năng nâng cao bao gồm nạp kịch bản mẫu để mô phỏng sự cố, xuất dữ liệu luồng ra bảng tính và phát hiện lỗi giao thức qua công cụ phân tích chuyên gia. Đồng thời, ứng dụng tự động khởi tạo luật tường lửa tương thích với thiết bị Cisco, iptables hoặc Windows Firewall, kết hợp mô hình học sâu FT-Transformer qua thư viện PyTorch để dự đoán hành vi luồng phục vụ công tác học tập.

\section{Tổng quan chức năng}
Dựa trên kết quả khảo sát hiện trạng, hệ thống được thiết kế để trở thành một phần mềm phân tích lưu lượng mạng đa năng. Mục tiêu cốt lõi là hỗ trợ đắc lực cho công tác thực hành kỹ thuật mạng, từ khâu thu thập dữ liệu thời gian thực, đọc phân tích tệp bắt gói, lọc thông tin theo cú pháp đến trực quan hóa đồ họa luồng và hỗ trợ dự đoán hành vi. Phần này cung cấp cái nhìn tổng thể về các khối chức năng mà hệ thống cung cấp cho người sử dụng.

\subsection{Biểu đồ ca sử dụng tổng quát}
\begin{figure}[H]
\centering
\includegraphics[width=0.85\textwidth]{Hinhve/All_Usecase.png}
\caption{Biểu đồ ca sử dụng tổng quát hệ thống}
\label{fig:usecase_tong_quat}
\end{figure}

Biểu đồ ca sử dụng tổng quát mô tả các chức năng mức cao của hệ thống phân tích gói tin được hiển thị tại Hình \ref{fig:usecase_tong_quat}.

\subsubsection{Các tác nhân hệ thống}
Hệ thống tương tác với tác nhân chính sau:
\begin{itemize}
\item \textbf{Người dùng:} Là đối tượng vận hành trực tiếp hệ thống, bao gồm sinh viên thực hành mạng, giảng viên giảng dạy lý thuyết, hoặc quản trị viên mạng cần giám sát hệ thống. Vai trò của người dùng là lựa chọn giao diện mạng, nạp tệp dữ liệu, thực hiện các tác vụ lọc cú pháp, theo dõi các thống kê trực quan và tham khảo dự đoán hành vi bất thường từ phân hệ học sâu.
\end{itemize}

\subsubsection{Các ca sử dụng chính}
Hệ thống phân tích gói tin cung cấp bảy ca sử dụng chính tương ứng với các phân hệ chức năng:
\begin{itemize}
\item \textbf{Bắt gói tin:} Cho phép quét chọn giao diện mạng cục bộ hoặc cấu hình qua giao thức SSH để thu thập gói tin từ xa. Người dùng có thể thiết lập bộ lọc bắt gói, chế độ nhận toàn bộ lưu lượng (Promiscuous Mode), và điều khiển phiên bắt gói thời gian thực.
\item \textbf{Thao tác tệp lưu trữ:} Hỗ trợ nạp, tải lại hoặc lưu trữ tệp tin pcap/pcapng và các kịch bản gói tin mẫu để giả lập sự cố. Chức năng này còn cho phép đóng phiên làm việc, ghép nhiều tệp tin hoặc phân tách tệp lớn theo dung lượng hoặc số lượng gói chỉ định.
\item \textbf{Phân tích gói tin:} Giải mã cấu trúc 13 giao thức mạng phổ biến theo mô hình phân lớp hiển thị chi tiết dạng cây và dạng dữ liệu dạng thập lục phân (hex) hoặc ký tự (ASCII) thô.
\item \textbf{Nhận diện gói tin:} Nhập và phân tích cú pháp bộ lọc hiển thị để sàng lọc danh sách gói tin động, kết hợp tìm kiếm chuỗi ký tự hoặc byte dữ liệu, hỗ trợ tự động tô màu dòng gói tin theo loại giao thức để tăng trực quan.
\item \textbf{Thống kê luồng giao tiếp:} Gom cụm gói tin thành các luồng giao tiếp logic hai chiều để tính toán thời gian, tổng số gói, dung lượng byte và tốc độ bit trung bình. Toàn bộ thông số đặc trưng luồng có thể được xuất ra CSV.
\item \textbf{Phát hiện hành vi bất thường:} Tự động phân tích các cảnh báo chuyên gia (Expert Info) để phát hiện lỗi giao thức hoặc bất thường truyền tải, kết hợp trích xuất đặc trưng luồng đưa vào mô hình học sâu FT-Transformer (PyTorch) để dự đoán hành vi bình thường hoặc các kiểu tấn công như từ chối dịch vụ (DoS) và quét cổng (PortScan).
\item \textbf{Trực quan hóa giao tiếp:} Cung cấp bảng điều khiển tổng quan trực quan bao gồm các biểu đồ dạng đường, cột, tròn thể hiện tỷ lệ giao thức và tốc độ lưu lượng. Đồng thời hỗ trợ vẽ lưu đồ tiến trình trao đổi dữ liệu (Flow Graph) và sơ đồ cấu trúc mạng liên kết (Topology Graph).
\end{itemize}

\subsection{Biểu đồ ca sử dụng phân rã bắt gói tin}
\begin{figure}[H]
\centering
\includegraphics[width=0.85\textwidth]{Hinhve/Usecase_BatGoi.png}
\caption{Biểu đồ ca sử dụng phân rã chức năng bắt gói tin}
\label{fig:usecase_bat_goi}
\end{figure}

Biểu đồ phân rã chức năng bắt gói tin tại Hình \ref{fig:usecase_bat_goi} mô tả chi tiết quy trình người dùng thiết lập và khởi chạy phiên thu thập lưu lượng mạng. Các ca sử dụng con phân rã bao gồm:
\begin{itemize}
\item \textbf{Chọn giao diện mạng:} Người dùng truy cập danh sách giao diện mạng vật lý hoặc ảo (như Wi-Fi, Ethernet, Loopback) được hệ thống tự động dò tìm để lựa chọn card mạng nguồn bắt gói tin trực tiếp.
\item \textbf{Kết nối giao diện từ xa:} Hỗ trợ người dùng nhập cấu hình thông tin kết nối bảo mật SSH (gồm IP, cổng, tài khoản, mật khẩu) và chỉ định giao diện mạng từ xa để thu thập gói tin từ các máy chủ hoặc thiết bị mạng ngoại vi khác.
\item \textbf{Chọn bộ lọc thích hợp:} Thiết lập các biểu thức bộ lọc bắt gói trước khi chạy giúp tối ưu hóa tài nguyên hệ thống bằng cách chỉ thu thập các gói tin thỏa mãn điều kiện mong muốn.
\item \textbf{Chọn cấu hình thích hợp:} Hỗ trợ cấu hình các thông số phần cứng và giao thức của card mạng như kích hoạt hoặc tắt chế độ nhận toàn bộ lưu lượng (Promiscuous Mode) và các thông số kích thước bộ đệm.
\item \textbf{Chọn chế độ vận hành:} Cấu hình các điều kiện dừng thu thập gói tin tự động (như giới hạn số lượng gói bắt được, giới hạn dung lượng tệp tin, hoặc giới hạn thời lượng bắt gói) và thiết lập ghi tệp tự động luân phiên (ring buffer) phù hợp với mục đích giám sát mạng dài hạn hoặc thực hành ngắn hạn.
\end{itemize}

\subsection{Biểu đồ ca sử dụng phân rã thao tác tệp lưu trữ}
\begin{figure}[H]
\centering
\includegraphics[width=0.85\textwidth]{Hinhve/Usecase_ThaoTacTep.png}
\caption{Biểu đồ ca sử dụng phân rã chức năng thao tác tệp lưu trữ}
\label{fig:usecase_thao_tac_tep}
\end{figure}

Biểu đồ phân rã chức năng thao tác tệp lưu trữ tại Hình \ref{fig:usecase_thao_tac_tep} mô tả chi tiết các tác vụ của người dùng khi làm việc với tệp tin dữ liệu mạng tĩnh. Các ca sử dụng con phân rã bao gồm:
\begin{itemize}
\item \textbf{Mở tệp:} Cho phép người dùng tìm kiếm và nạp các tệp tin dữ liệu lưu lượng mạng đã lưu từ trước dưới các định dạng chuẩn (như tệp pcap hoặc pcapng) vào hệ thống để bắt đầu quá trình giải mã, phân tích hoặc tải kịch bản gói tin mẫu giả lập sự cố.
\item \textbf{Lưu tệp:} Cung cấp hai lựa chọn bao gồm lưu đè trực tiếp lên tệp hiện tại (Save) để lưu trữ nhanh các gói tin mới bắt được, hoặc lưu thành một tệp tin mới với đường dẫn khác (Save As) để tránh làm thay đổi tệp tin gốc.
\item \textbf{Chọn định dạng tệp khi lưu:} Hỗ trợ người dùng lựa chọn đuôi mở rộng khi xuất tệp tin, kết hợp tùy chọn các chế độ nén dữ liệu như gzip, lz4 để tối ưu dung lượng lưu trữ trên đĩa cứng.
\item \textbf{Tách tệp:} Hỗ trợ phân rã một tệp dữ liệu bắt gói tin dung lượng lớn thành nhiều tệp nhỏ hơn dựa trên các tiêu chí cấu hình như giới hạn số lượng gói tin trên mỗi tệp hoặc giới hạn kích thước dung lượng tệp cụ thể.
\item \textbf{Ghép tệp:} Cho phép người dùng lựa chọn và kết hợp nội dung của nhiều tệp tin pcap hoặc pcapng riêng lẻ thành một tệp tin hợp nhất duy nhất theo thứ tự thời gian của các gói tin để thuận tiện cho việc giám sát toàn cảnh.
\end{itemize}

\subsection{Biểu đồ ca sử dụng phân rã phân tích gói tin}
\begin{figure}[H]
\centering
\includegraphics[width=0.85\textwidth]{Hinhve/Usecase_PhanTichGoi.png}
\caption{Biểu đồ ca sử dụng phân rã chức năng phân tích gói tin}
\label{fig:usecase_phan_tich_goi}
\end{figure}

Biểu đồ phân rã chức năng phân tích gói tin tại Hình \ref{fig:usecase_phan_tich_goi} thể hiện quy trình bóc tách và giải mã dữ liệu mạng. Các ca sử dụng con phân rã bao gồm:
\begin{itemize}
\item \textbf{Phân tích chi tiết:} Cho phép bóc tách sâu cấu trúc các trường thông tin của gói tin được chọn theo mô hình phân lớp giao thức (bao gồm Ethernet, IPv4, IPv6, TCP, UDP, ...) và hiển thị dưới dạng cây thư mục cấu trúc logic phân tầng để người dùng dễ dàng quan sát giá trị từng trường dữ liệu.
\item \textbf{Ánh xạ thông tin với byte:} Đồng bộ hóa dữ liệu giữa các trường thông tin trong cây phân tích với vùng dữ liệu thô dạng thập lục phân hexa và ký tự ASCII; khi người dùng nhấn chọn một trường cụ thể, hệ thống sẽ tự động tô đậm vùng byte tương ứng trong tệp dữ liệu thô giúp hiểu rõ cơ chế đóng gói.
\item \textbf{Xem danh sách gói tin:} Hiển thị bảng danh sách toàn bộ các gói tin đã bắt được theo thời gian thực hoặc nạp từ tệp tin tĩnh với các cột thông số cơ bản như số thứ tự gói, nhãn thời gian, địa chỉ IP nguồn hoặc đích, giao thức mạng, dung lượng byte và mô tả tóm tắt thông tin gói tin.
\end{itemize}

\subsection{Biểu đồ ca sử dụng phân rã nhận diện gói tin}
\begin{figure}[H]
\centering
\includegraphics[width=0.85\textwidth]{Hinhve/Usecase_LocTimKiem.png}
\caption{Biểu đồ ca sử dụng phân rã chức năng nhận diện gói tin}
\label{fig:usecase_loc_tim_kiem}
\end{figure}

Biểu đồ phân rã chức năng nhận diện gói tin tại Hình \ref{fig:usecase_loc_tim_kiem} mô tả chi tiết các tác vụ của người dùng nhằm tìm kiếm, sàng lọc và đánh dấu lưu lượng mạng cần chú ý. Các ca sử dụng con phân rã bao gồm:
\begin{itemize}
\item \textbf{Lọc gói tin:} Hỗ trợ nhập và kiểm tra cú pháp biểu thức lọc trên thanh công cụ để lọc danh sách gói tin thời gian thực hoặc tệp tin tĩnh. Hệ thống hỗ trợ lưu trữ danh sách các bộ lọc yêu thích đã thiết lập trước đó để gọi lại nhanh chóng khi cần.
\item \textbf{Tìm kiếm gói tin:} Cung cấp nhiều chế độ tìm kiếm bao gồm tìm kiếm theo chuỗi ký tự, tìm kiếm theo chuỗi thập lục phân hoặc tìm các byte cụ thể bên trong vùng dữ liệu thô. Người dùng có thể tìm kiếm xuôi hoặc ngược trong danh sách gói tin và tùy chọn phân biệt chữ hoa chữ thường.
\item \textbf{Di chuyển nhanh:} Cho phép chuyển hướng định vị nhanh chóng bằng cách nhập số thứ tự gói tin, nhảy về gói đầu tiên hoặc cuối cùng. Ngoài ra, hệ thống còn hỗ trợ di chuyển theo lịch sử xem trước đó, nhảy đến các gói tin tương ứng trong cùng một phiên hội thoại truyền tải hoặc gói tin yêu cầu/phản hồi liên quan.
\item \textbf{Đánh dấu và ẩn gói:} Hỗ trợ đánh dấu các gói tin quan trọng bằng cờ màu, tạm thời bỏ qua hoặc ẩn các gói tin nhiễu để thu gọn không gian hiển thị danh sách. Hệ thống cũng cho phép người dùng đính kèm các ghi chú văn bản chi tiết trực tiếp vào metadata của gói tin được chọn hoặc xóa nhanh toàn bộ ghi chú trong phiên làm việc.
\item \textbf{Tô màu gói tin:} Cung cấp hệ thống quy tắc tô màu tự động dựa trên giao thức của gói tin. Người dùng có thể chỉnh sửa các bảng màu mặc định, thêm mới hoặc thay đổi độ ưu tiên của các quy tắc tô màu tùy chỉnh để phân loại trực quan nhanh danh sách hiển thị.
\end{itemize}

\subsection{Biểu đồ ca sử dụng phân rã thống kê luồng giao tiếp}
\begin{figure}[H]
\centering
\includegraphics[width=0.85\textwidth]{Hinhve/Usecase_ThongKeLuong.png}
\caption{Biểu đồ ca sử dụng phân rã chức năng thống kê luồng giao tiếp}
\label{fig:usecase_thong_ke_luong}
\end{figure}

Biểu đồ phân rã chức năng thống kê luồng giao tiếp tại Hình \ref{fig:usecase_thong_ke_luong} mô tả chi tiết các tác vụ của người dùng nhằm phân tích các chỉ số truyền thông trên toàn hệ thống. Các ca sử dụng con phân rã bao gồm:
\begin{itemize}
\item \textbf{Thống kê hội thoại:} Gom nhóm các gói tin theo từng phiên truyền thông hai chiều giữa hai nút mạng cụ thể để tổng hợp số lượng gói tin, tổng số byte, tốc độ truyền tải cùng hướng truyền cho từng phiên.
\item \textbf{Thống kê đầu cuối:} Liệt kê danh sách tất cả các điểm kết nối mạng đơn lẻ bao gồm địa chỉ IP hoặc địa chỉ vật lý tham gia gửi hoặc nhận dữ liệu mạng, kèm theo tổng dung lượng byte và số lượng gói tin tương ứng của từng điểm.
\item \textbf{Thống kê địa chỉ và dịch vụ:} Phân tích và hiển thị thông tin thống kê về các địa chỉ IPv4, địa chỉ IPv6, danh sách các tên miền dịch vụ được phân giải cùng các dịch vụ giao thức lớp ứng dụng như HTTP.
\item \textbf{Thống kê phân cấp giao thức:} Phân tích tỷ lệ phần trăm phân bố của các giao thức mạng trong toàn bộ lưu lượng thu thập được dưới dạng cấu trúc cây phân lớp từ lớp liên kết dữ liệu đến lớp ứng dụng.
\item \textbf{Thống kê độ dài và chiều truyền:} Thống kê sự phân bố kích thước độ dài của các gói tin theo các khoảng kích cỡ được định nghĩa trước và phân tích tỷ lệ lượng dữ liệu truyền đi hoặc nhận về trên các kênh truyền thông.
\end{itemize}

\subsection{Biểu đồ ca sử dụng phân rã phát hiện hành vi bất thường}
\begin{figure}[H]
\centering
\includegraphics[width=0.85\textwidth]{Hinhve/Usecase_PhatHienHanhVi.png}
\caption{Biểu đồ ca sử dụng phân rã chức năng phát hiện hành vi bất thường}
\label{fig:usecase_phat_hien_hanh_vi}
\end{figure}

Biểu đồ phân rã chức năng phát hiện hành vi bất thường tại Hình \ref{fig:usecase_phat_hien_hanh_vi} mô tả các tác vụ liên quan đến phân tích thông minh và phát hiện sự cố hệ thống. Các ca sử dụng con phân rã bao gồm:
\begin{itemize}
\item \textbf{Định nghĩa luật phát hiện hành vi:} Cho phép người dùng thiết lập, chỉnh sửa, thêm mới hoặc xóa bỏ các điều kiện lọc, nhãn thông điệp cảnh báo và mức độ nghiêm trọng trong bảng luật chuyên gia.
\item \textbf{Phát hiện hành vi qua AI:} Sử dụng dữ liệu đặc trưng luồng mạng được trích xuất để nạp vào mô hình học sâu FT-Transformer huấn luyện bằng PyTorch, thực hiện dự đoán phân loại hành vi luồng thành các nhãn bình thường hoặc tấn công từ chối dịch vụ và quét cổng.
\item \textbf{Phát hiện hành vi qua tập luật:} Tự động đối chiếu các trường dữ liệu của gói tin với danh sách các điều kiện trong tập luật chuyên gia đã định nghĩa để phát hiện lỗi giao thức và đưa ra các cảnh báo lỗi truyền dẫn tương ứng.
\end{itemize}

\subsection{Biểu đồ ca sử dụng phân rã trực quan hóa giao tiếp}
\begin{figure}[H]
\centering
\includegraphics[width=0.85\textwidth]{Hinhve/Usecase_TrucQuanHoa.png}
\caption{Biểu đồ ca sử dụng phân rã chức năng trực quan hóa giao tiếp}
\label{fig:usecase_truc_quan_hoa}
\end{figure}

Biểu đồ phân rã chức năng trực quan hóa giao tiếp tại Hình \ref{fig:usecase_truc_quan_hoa} mô tả chi tiết các tác vụ đồ họa hóa thông tin trao đổi dữ liệu. Các ca sử dụng con phân rã bao gồm:
\begin{itemize}
\item \textbf{Trực quan hóa hội thoại theo thời gian:} Vẽ lưu đồ tiến trình trao đổi dữ liệu biểu diễn chuỗi thời gian gửi nhận gói tin tuần tự giữa các địa chỉ IP của các phiên hội thoại.
\item \textbf{Trực quan hóa quan hệ giữa các đầu cuối:} Dựng sơ đồ liên kết các nút thiết bị mạng dưới dạng đồ thị nút và các cạnh kết nối thể hiện mối liên hệ logic truyền thông giữa các máy.
\item \textbf{Điều chỉnh bố cục và chế độ hiển thị:} Cho phép người dùng chuyển đổi các thuật toán sắp xếp vị trí nút mạng, thu phóng sơ đồ kết nối, lọc hiển thị nhãn địa chỉ hoặc thay đổi định dạng biểu đồ thống kê.
\item \textbf{Trực quan bằng giao diện dashboard:} Tổng hợp các biểu đồ dạng đường thể hiện tốc độ truyền tải gói tin theo thời gian thực, biểu đồ cột biểu diễn lượng dữ liệu qua các cổng dịch vụ và biểu đồ tròn thống kê tỷ lệ phần trăm các giao thức lớp ứng dụng.
\end{itemize}

\subsection{Quy trình nghiệp vụ phân tích lưu lượng mạng}
Bên cạnh việc phân tách các ca sử dụng đơn lẻ, hệ thống còn vận hành dựa trên các quy trình nghiệp vụ phối hợp đa tác nhân. Quy trình quan trọng nhất là quy trình Phân tích lưu lượng mạng, mô tả vòng đời của một phiên giám sát mạng từ lúc khởi tạo nguồn dữ liệu đến khi kết thúc phân tích và xuất kết quả.

\begin{figure}[H]
\centering
\includegraphics[width=0.55\textwidth]{Hinhve/Activity_QuyTrinhNghiepVu.png}
\caption{Biểu đồ hoạt động mô tả quy trình nghiệp vụ phân tích lưu lượng mạng}
\label{fig:activity_quy_trinh}
\end{figure}

Quy trình bắt đầu bằng bước thiết lập nguồn dữ liệu, trong đó người dùng tiến hành dò tìm và lựa chọn giao diện mạng cục bộ phù hợp hoặc cấu hình các tham số bảo mật để kết nối với giao diện mạng từ xa. Đồng thời, người dùng có thể thiết lập thêm các biểu thức bộ lọc bắt gói để đảm bảo hệ thống chỉ thu thập đúng các lưu lượng mạng cần phân tích, tránh quá tải tài nguyên.

Tiếp theo là bước thu thập và phân tích dữ liệu, lúc này hệ thống sẽ kích hoạt phiên bắt gói tin trực tiếp thời gian thực hoặc nạp dữ liệu từ các tệp tin tĩnh có sẵn. Hệ thống tự động bóc tách và giải mã các gói tin thô thành các trường thông tin rõ ràng theo phân lớp giao thức mạng để người dùng có thể xem cấu trúc chi tiết dạng cây hoặc dạng dữ liệu thô.

Sau đó, hệ thống chuyển sang bước tổng hợp quan hệ và phân tích luồng bằng cách gom cụm các gói tin đơn lẻ thành các luồng giao tiếp hai chiều. Từ các luồng này, hệ thống tính toán các chỉ số truyền thông đầu cuối, phân cấp giao thức và đưa dữ liệu đặc trưng qua phân hệ phát hiện hành vi bất thường để kiểm tra tập luật hoặc dùng mô hình học sâu FT-Transformer dự đoán các hành vi tấn công.

Cuối cùng là bước trực quan hóa và xuất kết quả, toàn bộ thông tin thống kê luồng và cảnh báo bất thường được đồng bộ hóa để hiển thị trực quan lên giao diện dashboard, lưu đồ tiến trình trao đổi dữ liệu và sơ đồ cấu trúc mạng liên kết. Người dùng kết thúc quy trình bằng việc xuất dữ liệu đặc trưng luồng ra tệp tin bảng tính hoặc lưu trữ các báo cáo kết quả phân tích sự cố.

\section{Đặc tả chức năng}
Phần này mô tả ngắn gọn các chức năng quan trọng nhất của hệ thống. Mỗi chức năng được trình bày tập trung vào mục đích, điều kiện sử dụng, cách thực hiện và kết quả sau cùng.

\subsection{Đặc tả chức năng bắt gói tin}
\begin{itemize}
\item Mục đích: cho phép người dùng thu thập lưu lượng mạng trực tiếp từ các giao diện mạng cục bộ hoặc từ xa.
\item Tác nhân: người dùng.
\item Điều kiện đầu vào: thiết bị có các giao diện mạng hoạt động.
\item Luồng chính: người dùng chọn card mạng, bắt đầu quá trình bắt gói, theo dõi danh sách gói tin hiển thị thời gian thực và dừng khi cần.
\item Kết quả: hệ thống hiển thị danh sách các gói tin thu được lên màn hình.
\end{itemize}

\subsection{Đặc tả chức năng thao tác tệp lưu trữ}
\begin{itemize}
\item Mục đích: quản lý lưu giữ dữ liệu mạng thông qua tệp tin hoặc kịch bản mẫu.
\item Tác nhân: người dùng.
\item Điều kiện đầu vào: có sẵn tệp tin pcap hoặc pcapng hợp lệ, hoặc chọn kịch bản mẫu có sẵn.
\item Luồng chính: người dùng mở tệp tin từ đĩa hoặc chọn nạp một kịch bản demo mẫu, hệ thống thực hiện đọc và tải dữ liệu lên giao diện. Người dùng cũng có thể lưu danh sách gói tin đã bắt thành tệp tin mới.
\item Kết quả: dữ liệu được tải lên hoặc được lưu trữ thành công trên đĩa.
\end{itemize}

\subsection{Đặc tả chức năng phân tích gói tin}
\begin{itemize}
\item Mục đích: giải mã chi tiết các trường thông tin giao thức và cảnh báo lỗi truyền dẫn.
\item Tác nhân: người dùng.
\item Điều kiện đầu vào: đã có danh sách gói tin được nạp hoặc bắt trực tiếp.
\item Luồng chính: người dùng chọn một gói tin cụ thể trên danh sách, hệ thống phân tích và hiển thị cấu trúc chi tiết các lớp giao thức, biểu diễn mã hex tương ứng và hiển thị các cảnh báo chuyên gia về lỗi giao thức nếu có.
\item Kết quả: người dùng nắm được nội dung chi tiết và các thông tin kỹ thuật của gói tin.
\end{itemize}

\subsection{Đặc tả chức năng nhận diện gói tin}
\begin{itemize}
\item Mục đích: sàng lọc nhanh danh sách gói tin theo nhu cầu quan sát.
\item Tác nhân: người dùng.
\item Điều kiện đầu vào: danh sách gói tin đang hiển thị trên giao diện.
\item Luồng chính: người dùng nhập biểu thức điều kiện vào ô tìm kiếm theo cú pháp bộ lọc hiển thị, hệ thống thực hiện phân tích cú pháp và lọc danh sách gói tin thỏa mãn điều kiện.
\item Kết quả: danh sách gói tin hiển thị được rút gọn theo đúng điều kiện tìm kiếm.
\end{itemize}

\subsection{Đặc tả chức năng thống kê luồng giao tiếp}
\begin{itemize}
\item Mục đích: gom nhóm các gói tin thành luồng hai chiều và tính các tham số thống kê.
\item Tác nhân: người dùng.
\item Điều kiện đầu vào: đã có dữ liệu gói tin trong hệ thống.
\item Luồng chính: hệ thống gom các gói tin thành luồng theo địa chỉ, cổng và giao thức, trích xuất đặc trưng và tính toán các tham số về dung lượng, tốc độ và thời gian của luồng mạng.
\item Kết quả: hiển thị bảng danh sách các luồng giao tiếp kèm theo số liệu thống kê chi tiết.
\end{itemize}

\subsection{Đặc tả chức năng phát hiện hành vi bất thường}
\begin{itemize}
\item Mục đích: gán nhãn dự đoán hành vi luồng lưu lượng mạng bằng học sâu.
\item Tác nhân: người dùng và mô hình học sâu.
\item Điều kiện đầu vào: dữ liệu đặc trưng luồng mạng đã được kết xuất trong hệ thống.
\item Luồng chính: hệ thống chuyển dữ liệu đặc trưng qua phân hệ học sâu sử dụng mô hình FT-Transformer, thực hiện dự đoán nhãn hành vi của luồng và xuất báo cáo thống kê kết quả.
\item Kết quả: hiển thị báo cáo thống kê kết quả phân loại hành vi luồng phục vụ quá trình phân tích.
\end{itemize}

\subsection{Đặc tả chức năng trực quan hóa giao tiếp}
\begin{itemize}
\item Mục đích: biểu diễn lưu lượng mạng trực quan qua biểu đồ và sơ đồ topo.
\item Tác nhân: người dùng.
\item Điều kiện đầu vào: dữ liệu mạng đã được thu thập hoặc nạp vào hệ thống.
\item Luồng chính: hệ thống xử lý số liệu để vẽ các biểu đồ đường, cột, tròn trên bảng điều khiển, vẽ sơ đồ tiến trình dòng chảy dữ liệu giữa các máy và dựng sơ đồ topo liên kết các thiết bị mạng.
\item Kết quả: giao diện hiển thị trực quan các đồ thị và sơ đồ kết nối mạng.
\end{itemize}

\section{Yêu cầu phi chức năng}
Ngoài các chức năng chính, hệ thống còn cần đáp ứng một số yêu cầu phi chức năng để có thể sử dụng ổn định và thuận tiện hơn trong thực tế.
\begin{itemize}
\item Dễ dùng: giao diện phải rõ ràng, bố cục hợp lý, người mới vẫn có thể làm quen nhanh.
\item Ổn định: khi mở tệp lớn hoặc bắt gói lâu, phần mềm vẫn cần hoạt động ổn định và không bị treo.
\item Nhanh: việc lọc, tìm kiếm và chuyển giữa các màn hình nên đủ mượt để người dùng không bị gián đoạn.
\item Dễ bảo trì: phần mềm cần được chia tách hợp lý để dễ sửa, dễ thêm chức năng và dễ kiểm tra lỗi.
\item Tương thích: ứng dụng phải hoạt động ổn định trên môi trường triển khai mà đồ án hướng đến.
\item An toàn dữ liệu: dữ liệu đã mở hoặc đã lưu cần được giữ nguyên trạng, tránh mất mát ngoài ý muốn.
\item Trình bày rõ: chữ, bảng và hình phải đủ lớn để người đọc có thể hiểu ngay khi xem báo cáo.
\end{itemize}

\section{Kết chương}
Chương này đã khảo sát hiện trạng giám sát và phân tích lưu lượng mạng, chỉ ra các khó khăn thường gặp khi làm việc với dữ liệu gói tin, đồng thời xác định các nhóm chức năng chính của phần mềm. Từ đó, đồ án hình thành được một hướng thiết kế rõ ràng hơn: vừa thu thập được dữ liệu, vừa cho phép xem lại, vừa hỗ trợ lọc, thống kê, trực quan hóa và nhận diện hành vi.

Các nội dung trong chương này là cơ sở để triển khai các chương sau, đặc biệt là phần nền tảng công nghệ, phần thiết kế hệ thống và phần đánh giá kết quả. Khi yêu cầu đã được làm rõ, việc xây dựng phần mềm sẽ nhất quán hơn và dễ kiểm tra hơn trong thực tế.

\end{document}
